\honest{Ends: An Introduction}{Rob Bensinger}{2015}

Value theory is the study of what people care about: our goals, tastes, pleasures, fears and ambitions, what we wish we cared about, and everyday values like art, food, sex and friendship; going to the movies with your friend Sam can be a value without being a moral one. \nb{Philosophers use ``value theory'' for the theory of what is good; the \textsc{Stanford Encyclopedia of Philosophy} gives three senses, and none is defined by what people care about. The definition is the book's own.}

We study our values because we do not act as we wish: our preferences conflict, we want different desires, and we lack the will, attention or insight to act as we would like. Humans are not consistent enough to have utility functions. \nb{The book gives the evidence later, in the posts on the Allais paradox.} We also misjudge what we actually wish: there is a gulf between how we think we wish to act and how we actually wish to. Philosophers disagree about what we want, what we ought to want, and even what ``ought'' means, and the history of moral theory is ``piled high with the corpses of failed Guiding Principles.'' That track record bodes ill for anyone who needs a reliable, usable account of their goals, whether to design safe AI, build institutions, choose a charity or decide which virtues to cultivate. \nb{In the 2009 PhilPapers survey, none of the three main moral theories drew more than about a quarter of philosophers. A majority leaned toward moral realism.}

The book has three sequences: ``Fake Preferences,'' ``Value Theory'' and ``Quantified Humanism.'' The last, on applying theory to practice, is the most important, because ``The cash value of a normative theory is how well it translates into normative practice.'' \nb{It is close to the pragmatist view of theories. William James used ``cash-value'' in this sense.} A deeper understanding of your values should make you better at fulfilling them, and at the least should not get in the way: ``What good would it be, then, to know what's good?'' When to trust snap judgments over theory has no catch-all answer; we will have to look at examples and learn the warning signs of where sophisticated theories fail and where naive feelings fail.

The first recurring question is where we should go. Yudkowsky coined ``fun theory'' for it: what our ideal future would look like, not just its government or moral code but the adventures we would go on and the music we would compose. It meets transhumanism, which I define as ``the view that we can radically improve the human condition if we make enough scientific and social progress,'' and debates such as hedonism, the pursuit of pleasure, against eudaimonia, general well-being; the book also discusses cryonics, mind uploading and space colonization. \nb{The standard definition, in the Transhumanist FAQ, adds desirability and names technologies that change human beings. Philosophers usually set hedonism against desire theories and objective-list theories, and count eudaimonist accounts among the latter.}

Utopia-planning has become ``rather passe,'' because it seems naive and because we are bad at building utopias; the word itself means ``non-place.'' But if we give up on utopia, ``it's not obvious'' that our short-term pursuits will add up to a good future. Value is not inevitable; creating and preserving it takes work. \nb{A well-known case for the other side is Karl Popper's preference for piecemeal reform aimed at present ills over a plan for the good society.}

The second question is how ends relate to means. In a game the journey matters, and we do not want to skip ahead to being declared the winner; in saving a family member's life the outcome does, and you would take a less enjoyable strategy that worked better. Often our values lie in how the world turns out, especially the parts that can love and hurt and want. Last, I ask what we build into ``valuable'' at the start.

A footnote gives a transhumanist argument: we could abolish ageing and disease within decades or centuries, which would put us in the position of organisms with negligible senescence, such as lobsters and Aldabra giant tortoises, so we should invest in it. It also describes Sandel's objection to enhancement to enhancement as the worry that it would make life feel like less of a ``gift.'' \nb{Sandel's own name for it is ``the hubris objection.'' Sandel's concern is the ``drive to mastery'' that enhancement expresses, which ``may even destroy'' an appreciation of human gifts.}
